% Documento final del proyecto integrador AgentProof.
% Plantilla: guía de trabajo final de grado de la USFQ (pregrado), versión LaTeX.
% El preámbulo se conserva como en la plantilla, salvo el estilo bibliográfico (IEEE).
% Compilar desde esta carpeta con:  latexmk   (ver .latexmkrc)
\documentclass[12pt,a4paper]{article}

\usepackage{fontspec}
\setmainfont{Times New Roman}
\usepackage[spanish,es-noindentfirst,es-nolists,es-noshorthands,es-tabla]{babel}
\usepackage{geometry}
\usepackage{setspace}
\usepackage{xcolor}
\usepackage{enumitem}
\usepackage{array}
\usepackage{longtable}
\usepackage{needspace}
\usepackage{graphicx}
\usepackage{float}   % permite [H]: la tabla o figura se coloca exactamente donde se escribe
% Figuras dibujadas como código, para versionarlas como texto
\usepackage{tikz}
\usetikzlibrary{positioning,arrows.meta}
% ---------- Índices automáticos (tabla de contenido, índice de tablas e índice de figuras) ----------
%   Los títulos del cuerpo se crean con \tituloapa, \subtituloapa y \subsubtituloapa
%   (niveles 1, 2 y 3 de la guía USFQ) y entran solos en la tabla de contenido.
%   Las tablas y figuras entran en sus índices al usar \caption{} dentro de table/figure.
%   Compilar dos veces (latexmk lo hace solo) para que los números de página se actualicen.
\usepackage[titles]{tocloft}
\renewcommand{\cftdotsep}{1.2}
\setlength{\cftbeforesecskip}{6pt}        % 6 pt antes de cada entrada (estilo "TDC 1" del .docx)
\setlength{\cftsecindent}{0pt}
\setlength{\cftsecnumwidth}{0pt}
\renewcommand{\cftsecfont}{\normalfont}
\renewcommand{\cftsecpagefont}{\normalfont}
\renewcommand{\cftsecleader}{\cftdotfill{\cftdotsep}}
\setlength{\cftbeforesubsecskip}{6pt}
\setlength{\cftsubsecindent}{0.5in}
\setlength{\cftsubsecnumwidth}{0pt}
\renewcommand{\cftsubsecfont}{\normalfont}
\renewcommand{\cftsubsecleader}{\cftdotfill{\cftdotsep}}
\setlength{\cftbeforesubsubsecskip}{6pt}
\setlength{\cftsubsubsecindent}{1in}
\setlength{\cftsubsubsecnumwidth}{0pt}
\renewcommand{\cftsubsubsecfont}{\normalfont}
\renewcommand{\cftsubsubsecleader}{\cftdotfill{\cftdotsep}}
% Entradas "Tabla 1. Título ..... 12" y "Figura 1. Título ..... 12"
\setlength{\cftbeforetabskip}{6pt}
\setlength{\cfttabindent}{0pt}
\renewcommand{\cfttabpresnum}{Tabla~}
\renewcommand{\cfttabaftersnum}{.}
\setlength{\cfttabnumwidth}{0.75in}
\renewcommand{\cfttableader}{\cftdotfill{\cftdotsep}}
\setlength{\cftbeforefigskip}{6pt}
\setlength{\cftfigindent}{0pt}
\renewcommand{\cftfigpresnum}{Figura~}
\renewcommand{\cftfigaftersnum}{.}
\setlength{\cftfignumwidth}{0.8in}
\renewcommand{\cftfigleader}{\cftdotfill{\cftdotsep}}
% Índices sin título propio (el título lo pone \caratula, con el formato de la plantilla)
\makeatletter
\newcommand{\tablaDeContenido}{\begingroup\@starttoc{toc}\endgroup}
\newcommand{\indiceDeTablas}{\begingroup\@starttoc{lot}\endgroup}
\newcommand{\indiceDeFiguras}{\begingroup\@starttoc{lof}\endgroup}
\makeatother
% Leyendas de tablas y figuras: "Tabla 1. Título" / "Figura 1. Título"
\usepackage{caption}
\captionsetup{labelsep=period,justification=raggedright,singlelinecheck=false,font=normalsize,skip=6pt}
\addto\captionsspanish{\renewcommand{\tablename}{Tabla}\renewcommand{\figurename}{Figura}}
\setlength{\LTpre}{0pt}\setlength{\LTpost}{0pt}
\usepackage{fancyhdr}
\usepackage{csquotes}
% ---------- Bibliografía: biblatex con estilo IEEE (requiere compilar con biber) ----------
%   Compilar:  latexmk   (latexmk ejecuta biber solo)
%   Citas:     \cite{clave}       -> [1]
%              \cite{a,b}         -> [1], [2]
%              \cite[p.~12]{clave} -> [1, p. 12]
%   La lista de referencias sale numerada en orden de aparición.
%   Para agregar archivos .bib: colocarlos en la carpeta bib/ y registrarlos aquí.
\usepackage[backend=biber,style=ieee]{biblatex}
% Textos de las referencias en línea, como en la planificación: "consultado el ..." y "Disponible en: ..."
\DefineBibliographyStrings{spanish}{urlseen = {consultado el}, url = {Disponible en}}
% Los preprints de arXiv se registran como @article con journaltitle = {arXiv:<id>}, para que
% IEEE imprima «Título,» arXiv:<id>, <año>. Las páginas web (@online) llevan fecha de consulta.
\addbibresource{bib/referencias.bib}
% \addbibresource{bib/otro-archivo.bib}   % <- añade aquí más archivos .bib
\usepackage[hidelinks,hypertexnames=false]{hyperref}
\usepackage{xurl}\urlstyle{same}

% Word no divide palabras con guiones
\hyphenpenalty=10000 \exhyphenpenalty=10000
\tolerance=2000 \emergencystretch=2em

% ---------- Colores del documento Word ----------
\definecolor{rojo}{HTML}{FF0000}
\definecolor{gris}{HTML}{808080}
\definecolor{enlace}{HTML}{0070C0}
\newcommand{\rojo}[1]{\textcolor{rojo}{#1}}
\newcommand{\enlace}[1]{\href{#1}{\textcolor{enlace}{\underline{#1}}}}

% ---------- Interlineado (Word: Times New Roman 12 pt = 13.8 pt por línea) ----------
\setstretch{0.952}                 % sencillo  -> 13.8 pt
\newcommand{\dobleesp}{\setstretch{1.903}} % doble -> 27.6 pt
\setlength{\parindent}{0pt}
\setlength{\parskip}{0pt}

% ---------- Líneas en blanco (párrafos vacíos de Word) ----------
\newcommand{\blank}{\par\vspace{13.8pt}}          % 12 pt
\newcommand{\blankxiv}{\par\vspace{16.1pt}}       % 14 pt
\newcommand{\blankx}{\par\vspace{11.5pt}}         % 10 pt
\newcommand{\blankdoble}{\par\vspace{27.6pt}}     % 12 pt a doble espacio

% ---------- Tamaños de letra ----------
\newcommand{\xviii}{\fontsize{18pt}{20.7pt}\selectfont}
\newcommand{\xvi}{\fontsize{16pt}{18.4pt}\selectfont}
\newcommand{\xiv}{\fontsize{14pt}{16.1pt}\selectfont}
\newcommand{\xii}{\fontsize{12pt}{13.8pt}\selectfont}
\newcommand{\xionce}{\fontsize{11pt}{12.7pt}\selectfont}
\newcommand{\xps}{\fontsize{10pt}{11.5pt}\selectfont}

% ---------- Estilos de título ----------
% "Título 1 / APA Nivel 1": centrado, negrilla, mayúsculas, 12 pt, 12 pt antes, 24 pt después
%   Uso: \tituloapa[Texto para la tabla de contenido]{TEXTO DEL TÍTULO}
%   Si se omite el argumento opcional, en el índice aparece el mismo texto del título.
\makeatletter
\newcommand{\tituloapa}[2][\@nil]{%
  \par\vspace{12pt}\refstepcounter{section}%
  {\centering\bfseries\xii #2\par}%
  \def\@tmp{#1}\ifx\@tmp\@nnil\addcontentsline{toc}{section}{#2}\else\addcontentsline{toc}{section}{#1}\fi
  \vspace{24pt}}
\makeatother
% "APA Nivel 2": alineado a la izquierda, negrilla, 12 pt (entra como subsección en el índice)
\newcommand{\subtituloapa}[1]{\par\vspace{6pt}\refstepcounter{subsection}{\bfseries\xii #1\par}\addcontentsline{toc}{subsection}{#1}\nopagebreak}
% "APA Nivel 3": con sangría, negrilla, termina en punto; el texto continúa en la misma línea
\newcommand{\subsubtituloapa}[1]{\par\vspace{6pt}\refstepcounter{subsubsection}\addcontentsline{toc}{subsubsection}{#1}\hspace*{0.5in}{\bfseries\xii #1.}\ }
% "Carátula Tesis" 12 pt: centrado, negrilla, 12 pt antes, 18 pt después
\newcommand{\caratula}[1]{\par\vspace{12pt}{\centering\bfseries\xii #1\par}\vspace{18pt}}
% "Title": centrado, mayúsculas, 18 pt, negrilla, 12 pt antes y después
\newcommand{\titulogrande}[1]{\par\vspace{12pt}{\centering\bfseries\xviii #1\par}\vspace{12pt}}

% ---------- Tablas ----------
\setlength{\tabcolsep}{0.08in}
\renewcommand{\arraystretch}{1}
\newcommand{\celda}[1]{\vspace{-0.5\baselineskip}#1\par\vspace{-0.5\baselineskip}}

% ---------- Encabezado con número de página (arriba a la derecha) ----------
\pagestyle{fancy}
\fancyhf{}
\fancyhead[R]{\xii\thepage}
\renewcommand{\headrulewidth}{0pt}
\setlength{\headheight}{14pt}

% Lista de referencias: sangría francesa 0.5 in, 12 pt entre entradas, sin título propio
\setlength{\bibhang}{0.5in}
\setlength{\bibitemsep}{12pt}
\setlength{\bibparsep}{0pt}
\defbibheading{sinTitulo}{}

% ---------- Datos del documento (portada, calificación y derechos de autor) ----------
\newcommand{\tituloTrabajo}{AgentProof: inducción automática de escenarios multi-turno desde trazas de interacción y reejecución adaptativa con usuarios simulados para la detección de regresiones en agentes basados en LLM}
\newcommand{\nombreAutor}{Jarod Styph Tierra Medina}
\newcommand{\nombreColegio}{Ciencias e Ingenierías}
\newcommand{\nombreCarrera}{Ingeniería en Ciencias de la Computación}
\newcommand{\tituloObtenido}{\rojo{XXXX}}                          % por confirmar
\newcommand{\nombreTutor}{Ricardo Flores Moyano, \rojo{título académico}} % título por confirmar
\newcommand{\fechaDocumento}{27 de septiembre de 2026}

% ---------- Datos personales (no se publican) ----------
% thesis/privado.tex, ignorado por git, define \codigoEstudiante y \cedulaIdentidad.
% Sin ese archivo (por ejemplo, en el repositorio público) se muestran marcadores.
% Si se crea después de haber compilado, forzar la recompilación con: latexmk -g
\IfFileExists{privado.tex}{\input{privado.tex}}{%
  \newcommand{\codigoEstudiante}{\rojo{xxxxxxxx}}%
  \newcommand{\cedulaIdentidad}{\rojo{xxxxxxxxxx}}}

\begin{document}

% Márgenes de 1 pulgada en todo el documento (guía USFQ)
\newgeometry{top=1in,bottom=1in,left=1in,right=1in,headheight=14pt,headsep=22pt,footskip=0.5in}
\setcounter{page}{1}

% ---------------------------- PORTADA ----------------------------
\thispagestyle{empty}
\begin{center}
{\bfseries\xviii UNIVERSIDAD SAN FRANCISCO DE QUITO USFQ\par}
\blank\blank
{\bfseries\xiv Colegio de \nombreColegio\par}
\blank\blank\blank\blank\blank\blank\blank
{\bfseries\xvi \tituloTrabajo\par}
\blank
{\xps\textcolor{gris}{.}\par}
\blank\blank\blank\blank
{\bfseries\xviii \nombreAutor\par}
\blank
{\bfseries\xvi \nombreCarrera\par}
\blank\blank\blank
{\xiv Trabajo de fin de carrera presentado como requisito\par}
{\xiv para la obtención del título de\par}
{\xiv \tituloObtenido\par}
\blankx\blankx\blankx\blankx\blankx
Quito, \fechaDocumento\par
\end{center}

% ------------------------ HOJA DE CALIFICACIÓN ------------------------
\clearpage
\pagestyle{fancy}
\blank\blank
\titulogrande{UNIVERSIDAD SAN FRANCISCO DE QUITO USFQ}
\begin{center}
{\bfseries\xiv Colegio de \nombreColegio\par}
\blankxiv\blankxiv
{\bfseries\xiv HOJA DE CALIFICACIÓN\par}
{\bfseries\xiv DE TRABAJO DE FIN DE CARRERA\par}
\blankxiv\blankxiv
{\bfseries\xiv \tituloTrabajo\par}
\blankxiv
\vspace{18pt}{\bfseries\xviii \nombreAutor\par}\vspace{18pt}
\end{center}
\blank\blank\blank\blank\blank\blank\blank\blank
{\bfseries\makebox[4in][l]{Nombre del profesor, Título académico}\nombreTutor\par}
\blank\blank\blank\blank\blank\blank\blank\blank\blank
{\centering Quito, \fechaDocumento\par}

% ------------------------ DERECHOS DE AUTOR ------------------------
\clearpage
\blank
\titulogrande{© DERECHOS DE AUTOR}
\begingroup\dobleesp\setlength{\parindent}{0.5in}
\vspace{6pt}
Por medio del presente documento certifico que he leído todas las Políticas y Manuales de la Universidad San Francisco de Quito USFQ, incluyendo la Política de Propiedad Intelectual USFQ, y estoy de acuerdo con su contenido, por lo que los derechos de propiedad intelectual del presente trabajo quedan sujetos a lo dispuesto en esas Políticas.
\par\vspace{6pt}
Asimismo, autorizo a la USFQ para que realice la digitalización y publicación de este trabajo en el repositorio virtual, de conformidad a lo dispuesto en la Ley Orgánica de Educación Superior del Ecuador.
\par
\endgroup
\blank\blank\blank\blank\blank\blank\blank\blank
\makebox[2.35in][l]{Nombres y apellidos:}\rojo{xxxxxxx\ \ \ \ xxxxx\ \ \ xxxxx\ \ xxxxxxx}\par
\blank\blank
\makebox[2.35in][l]{Código:}\rojo{xxxxxxx}\par
\blank\blank
\makebox[2.35in][l]{Cédula de identidad:}\rojo{xxxxxxxxx}\par
\blank\blank
\makebox[2.6in][l]{Lugar y fecha:}\rojo{Ciudad}, \rojo{día} de \rojo{mes} de \rojo{año}\par


% Páginas preliminares y cuerpo: la numeración continúa desde 4 (guía USFQ)
\clearpage
\setcounter{page}{4}
% ------------------- ACLARACIÓN PARA PUBLICACIÓN -------------------
\titulogrande{ACLARACIÓN PARA PUBLICACIÓN}
\begingroup\dobleesp
\vspace{6pt}
\textbf{Nota:} El presente trabajo, en su totalidad o cualquiera de sus partes, no debe ser considerado como una publicación, incluso a pesar de estar disponible sin restricciones a través de un repositorio institucional. Esta declaración se alinea con las prácticas y recomendaciones presentadas por el Committee on Publication Ethics COPE descritas por Barbour et al. (2017) Discussion document on best practice for issues around theses publishing, disponible en \enlace{http://bit.ly/COPETheses}.
\par
\endgroup
\titulogrande{UNPUBLISHED DOCUMENT}
\begingroup\dobleesp
\vspace{6pt}
\textbf{Note:} The following capstone project is available through Universidad San Francisco de Quito USFQ institutional repository. Nonetheless, this project -- in whole or in part -- should not be considered a publication. This statement follows the recommendations presented by the Committee on Publication Ethics COPE described by Barbour et al. (2017) Discussion document on best practice for issues around theses publishing available on \enlace{http://bit.ly/COPETheses}.
\par
\endgroup

% ----------------------------- RESUMEN -----------------------------
\clearpage
\blank
\caratula{RESUMEN}
\begingroup\dobleesp\color{rojo}
\vspace{6pt}
En texto normal, debes presentar una descripción completa pero concisa de tu trabajo, que motive a potenciales lectores a revisarlo por completo. El resumen debe indicar claramente cuál es el asunto tratado en el trabajo, haciendo referencia a las motivaciones y enfoques utilizados para su desarrollo, los resultados más destacables, y las principales conclusiones que indiquen las implicaciones actuales y perspectivas futuras del asunto.
\par\vspace{6pt}
\textbf{Palabras clave:} Deben incluir entre 5 y 10 palabras claves que describan tu artículo.
\par
\endgroup

% ----------------------------- ABSTRACT -----------------------------
\clearpage
\blank
\caratula{ABSTRACT}
\begingroup\dobleesp\color{rojo}
\vspace{6pt}
En texto normal, debe ser una traducción precisa del resumen.
\par\vspace{6pt}
\textbf{Key words:} Presentar una traducción precisa de las palabras clave.
\par
\endgroup

% ------------------------ TABLA DE CONTENIDO ------------------------
% Se genera automáticamente con los títulos creados con \tituloapa, \subtituloapa
% y \subsubtituloapa. No editar a mano.
\clearpage
\blank
\caratula{TABLA DE CONTENIDO}
\begingroup
\tablaDeContenido
\endgroup

% ------------------------- ÍNDICE DE TABLAS -------------------------
% Se genera automáticamente a partir de los \caption{} de los entornos table.
\clearpage
\blank
\caratula{ÍNDICE DE TABLAS}
\begingroup
\indiceDeTablas
\endgroup

% ------------------------ ÍNDICE DE FIGURAS ------------------------
% Se genera automáticamente a partir de los \caption{} de los entornos figure.
\clearpage
\blank
\caratula{ÍNDICE DE FIGURAS}
\begingroup
\indiceDeFiguras
\endgroup


% Cuerpo: capítulos del sumario de contenidos de la planificación (§6)
% --------------------------- INTRODUCCIÓN ---------------------------
% Contenido (sumario de la planificación, §6): Contexto, problema, objetivos y contribuciones.
\clearpage
\blank
\tituloapa[Introducción]{INTRODUCCIÓN}
\begingroup\dobleesp
\vspace{6pt}
\par
\endgroup

% ------------------------- ESTADO DEL ARTE -------------------------
% Contenido (sumario de la planificación, §6): Agentes LLM, evaluación de agentes, simulación de usuarios, LLM-as-a-Judge,
% taxonomías de fallos, herramientas existentes y brecha identificada.
\clearpage
\blank
\tituloapa[Estado del arte]{ESTADO DEL ARTE}
\begingroup\dobleesp
\vspace{6pt}
Este capítulo presenta los conceptos en los que se apoya el trabajo, las herramientas y los trabajos académicos más cercanos a la propuesta, y la brecha que la motiva.
\par

\subtituloapa{Conceptos principales}
Un agente basado en modelos de lenguaje (LLM) es un sistema en el que un modelo interactúa con herramientas externas para completar tareas de varios pasos. La literatura de diseño distingue dos patrones: los \textit{workflows}, cuyo flujo de control está predefinido en código, y los agentes propiamente dichos, en los que el modelo dirige su propio proceso en un bucle con retroalimentación del entorno \cite{anthropic2024agents}. La distinción importa aquí en dos sentidos: el sistema que se evalúa es un agente, y el propio AgentProof se diseña como un workflow que delega la cognición en puntos acotados, como se detalla en el capítulo de la propuesta.
\par\vspace{6pt}
La evaluación de agentes se apoya en dos técnicas centrales \cite{yehudai2025survey}. La primera es la simulación de usuarios: un modelo de lenguaje interpreta al usuario y sostiene conversaciones de prueba con el agente, lo que permite ejercitar interacciones de varios turnos sin que una persona intervenga en cada ejecución. La segunda es el patrón \textit{LLM-as-a-Judge}, en el que un modelo califica ejecuciones contra criterios definidos. La confiabilidad de este evaluador automático no puede darse por supuesta y debe validarse contra jueces humanos \cite{yehudai2025survey}; la del usuario simulado se discute más adelante, junto con los trabajos que la han medido.
\par\vspace{6pt}
Para caracterizar cómo fallan los agentes se dispone de taxonomías específicas. La del AI Red Team de Microsoft cubre agentes individuales que usan herramientas \cite{msredteam2025taxonomy}. Trabajos recientes documentan además fallos silenciosos en sistemas en producción: errores cuya señal nunca llega a una persona de forma accionable y que, en su variante más peligrosa, el propio modelo convierte en una respuesta fluida y verosímil para el usuario \cite{wu2026narratives}. Estas taxonomías dan la semántica de las degradaciones controladas que se emplean en la evaluación experimental de este trabajo.
\par
\endgroup

% ------------------- DESCRIPCIÓN DE LA PROPUESTA -------------------
% Contenido (sumario de la planificación, §6): Arquitectura del pipeline, patrón de diseño híbrido, esquema de trazas,
% protocolo de sesión y adaptadores, selección de modelos.
\clearpage
\blank
\tituloapa[Descripción de la propuesta]{DESCRIPCIÓN DE LA PROPUESTA}
\begingroup\dobleesp
\vspace{6pt}
\par
\endgroup

% --------------------- DESARROLLO DEL PROTOTIPO ---------------------
% Contenido (sumario de la planificación, §6): Implementación de los módulos, agentes de caso de estudio, operadores de
% mutación y pruebas automatizadas.
\clearpage
\blank
\tituloapa[Desarrollo del prototipo]{DESARROLLO DEL PROTOTIPO}
\begingroup\dobleesp
\vspace{6pt}
\par
\endgroup

% ---------------- DISEÑO EXPERIMENTAL Y PRE-REGISTRO ----------------
% Contenido (sumario de la planificación, §6): Corpus semilla, conjuntos expuesto y manifestado, predicados de fallo,
% criterio de regresión y tamaño muestral.
\clearpage
\blank
\tituloapa[Diseño experimental y pre-registro]{DISEÑO EXPERIMENTAL Y PRE-REGISTRO}
\begingroup\dobleesp
\vspace{6pt}
\par
\endgroup

% -------------- EXPERIMENTOS Y ANÁLISIS DE RESULTADOS --------------
% Contenido (sumario de la planificación, §6): E1 a E4, con costos y latencias.
\clearpage
\blank
\tituloapa[Experimentos y análisis de resultados]{EXPERIMENTOS Y ANÁLISIS DE RESULTADOS}
\begingroup\dobleesp
\vspace{6pt}
\par
\endgroup

% ------------------ CONCLUSIONES Y TRABAJO FUTURO ------------------
% Contenido (sumario de la planificación, §6): Limitaciones declaradas, adaptadores adicionales, telemetría OpenTelemetry,
% ablación con autocrítica y validación sobre tráfico de producción.
\clearpage
\blank
\tituloapa[Conclusiones y trabajo futuro]{CONCLUSIONES Y TRABAJO FUTURO}
\begingroup\dobleesp
\vspace{6pt}
\par
\endgroup


% ---------------------- REFERENCIAS BIBLIOGRÁFICAS ----------------------
% La lista se genera sola con las obras citadas en el texto (bib/referencias.bib).
\clearpage
\blank
\tituloapa[Referencias bibliográficas]{REFERENCIAS BIBLIOGRÁFICAS}
\printbibliography[heading=sinTitulo]

\end{document}
